% Part: first-order-logic
% Chapter: syntax-and-semantics
% Section: first-order-languages

\documentclass[../../../include/open-logic-section]{subfiles}

\begin{document}

\olfileid{fol}{syn}{fol}

\olsection{一阶语言}

一阶逻辑的表达式由一个基本词汇构建而成，
该词汇包含\emph{变元}、\emph{常项符号}、
\emph{谓词符号}，有时还有\emph{函数符号}。
由它们，连同逻辑联结词、量词以及诸如括号和逗号的标点符号，
便构成了\emph{项}和\emph{公式}。

\begin{explain}
非形式地说，谓词符号是性质和关系的名称，
常项符号是个体对象的名称，而函数符号是映射的名称。
除了等词~$\eq$ 之外，这些符号都是\emph{非逻辑符号}，它们共同构成一个语言。
任何一个一阶语言~$\Lang L$ 都由其非逻辑符号决定。
在最一般的情况下，$\Lang L$ 的每类符号都包含无穷多个。
\end{explain}

在一般情况下，一阶逻辑中使用以下符号：

\begin{enumerate}
\item 逻辑符号
\begin{enumerate}
\item 逻辑联结词：
  \startycommalist
  \iftag{prvNot}{\ycomma $\lnot$（否定）}{}%
  \iftag{prvAnd}{\ycomma $\land$（合取）}{}%
  \iftag{prvOr}{\ycomma $\lor$（析取）}{}%
  \iftag{prvIf}{\ycomma $\lif$（条件句）}{}%
  \iftag{prvIff}{\ycomma $\liff$（双条件句）}{}%
  \iftag{prvAll}{\ycomma $\lforall$（全称量词）}{}%
  \iftag{prvEx}{\ycomma $\lexists$（存在量词）}{}.
\tagitem{prvFalse}{假命题常项~$\lfalse$。}{}
\tagitem{prvTrue}{真命题常项~$\ltrue$。}{}
\item 二元等词~$\eq$。
\item 一个可数无穷的变元集合：$\Obj v_0$、$\Obj v_1$、$\Obj v_2$、……。
\end{enumerate}
\item 非逻辑符号，构成一阶逻辑的\emph{标准语言}
\begin{enumerate}
\item 对每个 $n>0$，有一个可数无穷的 $n$元谓词符号集合：$\Obj A^n_0$、$\Obj A^n_1$、$\Obj A^n_2$、……。
\item 一个可数无穷的常项符号集合：$\Obj c_0$、$\Obj c_1$、$\Obj c_2$、……。
\item 对每个 $n>0$，有一个可数无穷的 $n$元函数符号集合：$\Obj f^n_0$、$\Obj f^n_1$、$\Obj f^n_2$、……。
\end{enumerate}
\item 标点符号：（、）、以及逗号。
\end{enumerate}

我们的大部分定义和结论都是针对一阶逻辑的完整标准语言给出的。不过，根据具体应用的需要，我们也可以将语言限制为只包含少数几个谓词符号、常项符号和函数符号。

\begin{ex}
算术语言~$\Lang L_A$ 包含一个二元谓词符号~$<$，一个常项符号~$\Obj 0$，一个一元函数符号~$\prime$，以及两个二元函数符号~$+$ 和~$\times$。
\end{ex}

\begin{ex}
集合论语言~$\Lang L_Z$ 只包含一个二元谓词符号~$\in$。
\end{ex}

\begin{ex}
序语言~$\Lang L_\le$ 只包含一个二元谓词符号~$\le$。
\end{ex}

同样地，这些只是约定：严格来说，它们不过是别名，例如，$<$、$\in$ 和~$\le$ 是~$\Obj A^2_0$ 的别名，$\Obj 0$ 是~$\Obj c_0$ 的别名，$\prime$ 是~$\Obj f^1_0$ 的别名，$+$ 是~$\Obj f^2_0$ 的别名，$\times$ 是~$\Obj f^2_1$ 的别名。

\iftag{defNot,defOr,defAnd,defIf,defIff,defTrue,defFalse,defEx,defAll}{%
除了上面引入的原始联结词
\iftag{notprvEx,notprvAll}{与量词}{与量词}，
我们还使用以下\emph{被定义}符号：
\startycommalist
  \iftag{defNot}{\ycomma $\lnot$ (否定)}{}%
  \iftag{defAnd}{\ycomma $\land$ (合取)}{}%
  \iftag{defOr}{\ycomma $\lor$ (析取)}{}%
  \iftag{defIf}{\ycomma $\lif$ (条件句)}{}%
  \iftag{defIff}{\ycomma $\liff$ (双条件句)}{}%
  \iftag{defAll}{\ycomma $\lforall$ (全称量词)}{}%
  \iftag{defEx}{\ycomma $\lexists$ (存在量词)}{}%
  \iftag{defFalse}{\ycomma 假~$\lfalse$}{}%
  \iftag{defTrue}{\ycomma 真~$\ltrue$}}{}.

\begin{tagblock}{defNot,defOr,defAnd,defIf,defIff,defTrue,defFalse,defEx,defAll}
\begin{explain}
被定义符号严格来说并非语言的一部分，而是作为一种非正式的缩写引入的：它允许我们缩写那些仅用原始符号书写会变得相当冗长的公式。这显然是一个优点。然而，更大的好处在于证明会变短。如果一个符号是原始符号，在证明中就需要单独处理。因此，原始符号越多，我们的证明就越长。
\end{explain}
\end{tagblock}

% Alternate symbols

\begin{intro}
大家可能会觉得，上面我们使用的术语和符号，与你们熟悉的不太一样。
逻辑学教材（和老师）通常会用 $\sim$、$\neg$ 或~! 表示“否定”；
会用 $\wedge$、$\cdot$ 或 $\&$ 表示“合取”。
常用的“条件句”或“蕴涵”符号有 $\rightarrow$、$\Rightarrow$ 和 $\supset$。
\iftag{prvIff,defIff}{“双条件句”“双蕴涵”或“（实质）等值”的符号有 $\leftrightarrow$、
  $\Leftrightarrow$ 和 $\equiv$。}{}
\iftag{prvFalse,defFalse}{符号 $\lfalse$ 常被称为“假”、“谬符”、“荒谬”或“底元”。}{}
\iftag{prvTrue,defTrue}{符号 $\ltrue$ 常被称为“真”、“真符”或“顶元”。}{}

按照惯例，用拉丁字母表开头的字母（例如 $a$、$b$、$c$）作为常项符号
（有时也称为名称），用拉丁字母表末尾的字母（例如 $x$、$y$、$z$）作为变元。
量词与变元结合，例如与 $x$ 结合；
记法上的变体包括：对于全称量词有 $\forall x$、$(\forall x)$、$(x)$、$\Pi x$、
$\bigwedge_x$；对于存在量词有 $\exists x$、$(\exists x)$、$(Ex)$、$\Sigma x$、
$\bigvee_x$。
\end{intro}

\begin{explain}
我们可以把所有命题联结词和两个量词都作为语言的原始符号。
我们也可以选择较少的一组原始符号，而将其他逻辑算子视为被定义的。
“真值函数完备的”布尔算子集合包括 $\{ \lnot, \lor \}$、$\{ \lnot, \land \}$
和 $\{ \lnot, \lif\}$——这些集合与任一量词结合，
都可以形成表达力完备的一阶语言。

大家可能还熟悉另外两个逻辑算子：Sheffer（谢弗）竖线~$|$（以 Henry Sheffer 的名字命名），
以及 Peirce（皮尔士）箭头~$\downarrow$，后者也被称为 Quine（蒯因）的匕首。
当它们按通常的读法分别作为“与非”和“或非”时，
这两个算子单凭自身就是真值函数完备的。
\end{explain}

\end{document}